
\documentclass[../../main.tex]{subfiles}
\begin{document}

% 年度の大きな表示
\begin{center}
  {\fontsize{48pt}{58pt}\selectfont \textbf{1972}\normalsize\textbf{年度}}
\end{center}

\vspace{10mm}

% 本文

\vspace{15mm}

% 出題テーマと難易度の表
\noindent\textbf{出題テーマと難易度}

\vspace{5mm}

\begin{center}
  \shadowbox{%
    \begin{tabular}{|c|p{8cm}|c|}
      \hline
      \textbf{問題番号} & \textbf{テーマ} & \textbf{難易度} \\
      \hline
      第1問 & 複素数と整数問題 & \\
      \hline
      第2問 & 複素数と方程式の解 & \\
      \hline
      第3問 & 二項係数と剰余の定理 & \\
      \hline
      第4問 & 積分関数の最小値 & \\
      \hline
      第5問 & 対数関数と回転体の体積 & \\
      \hline
      第6問 & 級数の収束条件と関数方程式 & \\
      \hline
    \end{tabular}%
  }
\end{center}

\vspace{15mm}

% 解答
\noindent\textbf{解答}

\vspace{5mm}

\begin{center}
  \shadowbox{%
    \renewcommand{\arraystretch}{1.6}
    \begin{tabular}{|c|c|p{8cm}|}
      \hline
      \textbf{問題番号} & \textbf{小問} & \textbf{答え} \\
      \hline
      第1問 & - & $(a,b)=(0,\pm1),(1,1),(1,0),(-1,-1),(-1,0)$ \\
      \hline
      第2問 & - & $x=-\sqrt{2},\ \dfrac{\sqrt{2}\pm\sqrt{2}i}{2}$ \\
      \hline
      \multirow{2}{*}{第3問} & (1) & $l>m$ \\
      \cline{2-3}
                             & (2) & $R(x)=\dfrac{1}{2}F''(a)(x-a)^2+F'(a)(x-a)+F(a)$ \\
      \hline
      第4問 & - & $x=\dfrac{a+b}{2}$ \\
      \hline
      第5問 & - & $V=\pi\left(2-\dfrac{2e}{3}\right)$ \\
      \hline
      \multirow{2}{*}{第6問} & (1) & $d\le0$ \\
      \cline{2-3}
                             & (2) & $f(x)=x-\pi$ \\
      \hline
    \end{tabular}%
    \renewcommand{\arraystretch}{1.0}
  }
\end{center}

\vspace{15mm}

% 解答時間
\noindent\textbf{試験形態}

\vspace{5mm}

\begin{center}
  \shadowbox{%
    \begin{tabular}{p{8cm}}
      （要確認）
    \end{tabular}%
  }
\end{center}

\end{document}
